% sections/supporting.tex -- captions of the supporting information (the items themselves are in supplement.tex / S1 Appendix).
% Numbered by type in order of first citation in the text; supplement.tex includes the items in the same order.
\paragraph*{S1 Appendix.}
\label{si:appendix}
{\bf Supplementary methods, tables and figures.} Procedures of the data audit, FOV mask, resizing comparison, loss check,
SAM implementation, seam check, oracle prompts, cost measurement and reproducibility, followed by all supporting tables
and figures listed here (supplement.pdf).

\paragraph*{S1 Checklist.}
\label{si:claim}
{\bf CLAIM 2024 checklist.} Location in the manuscript of every item of the Checklist for Artificial Intelligence in
Medical Imaging (2024 update).

\paragraph*{S1 Table.}
\label{si:resize}
{\bf Training-mask resizing.} Thirty training masks reduced by five methods, enlarged back to 2048 and scored against the
originals.

\paragraph*{S2 Table.}
\label{si:pilot}
{\bf Resolution pilot.} U-Net (seed 0) at 512 and 1024 on the validation set, with the decision rule fixed beforehand.

\paragraph*{S3 Table.}
\label{si:seam}
{\bf Tile-seam check.} Discontinuity at the tile borders and its effect on accuracy near the borders (validation, seed 0).

\paragraph*{S4 Table.}
\label{si:oracle}
{\bf Oracle-prompted zero-shot SAM and mask-resampling references.} Reported separately from the automatic models.

\paragraph*{S5 Table.}
\label{si:runs}
{\bf Training runs.} Epochs, best epoch, selected threshold, best validation Dice and validation thin-vessel recall of every
run.

\paragraph*{S6 Table.}
\label{si:overlap}
{\bf IoU, precision and recall on the test set.}

\paragraph*{S7 Table.}
\label{si:imagelevel}
{\bf Image-level spread of the test results.} Standard deviation and median over the test images.

\paragraph*{S8 Table.}
\label{si:width}
{\bf Width classes with alternative cut-offs (4 and 8 pixels).}

\paragraph*{S9 Table.}
\label{si:simple}
{\bf Simple effects in the 2 $\times$ 2 ablation.} Effect of each module with and without the other.

\paragraph*{S1 Fig.}
\label{si:training}
{\bf Validation Dice during training.} Every training run, per epoch.

\paragraph*{S10 Table.}
\label{si:fixed}
{\bf Test results with every run at a fixed threshold of 0.5.}

\paragraph*{S2 Fig.}
\label{si:thrval}
{\bf Thin-vessel recall against the threshold on the validation set.}

\paragraph*{S11 Table.}
\label{si:disease}
{\bf Test results by disease category.}

\paragraph*{S12 Table.}
\label{si:agreement}
{\bf Agreement between two available annotations of the same image.} Post hoc descriptive analysis using labels only
(n = \AgreeNPairs{}).

\paragraph*{S3 Fig.}
\label{si:qualall}
{\bf All qualitative test images.} Native-resolution crops and error maps of every model (seed 0).

\paragraph*{S4 Fig.}
\label{si:diseasefig}
{\bf Results by disease category (test set).}

\paragraph*{S13 Table.}
\label{si:vessel}
{\bf Image-derived vessel measures.}

\paragraph*{S5 Fig.}
\label{si:blandaltman}
{\bf Bland--Altman plots of vessel density and fractal dimension.}

\paragraph*{S6 Fig.}
\label{si:costfig}
{\bf Inference cost and thin-vessel recall.}

\paragraph*{S14 Table.}
\label{si:matched}
{\bf Thin-vessel recall at matched precision.} Post hoc analysis; thresholds matched on the validation set.

\paragraph*{S7 Fig.}
\label{si:precthin}
{\bf Thin-vessel recall against overall precision (test set).} Post hoc analysis.

\paragraph*{S15 Table.}
\label{si:precision}
{\bf Width-stratified precision, spur-pruned thin class, and clDice and NSD at 0.5.} Post hoc analysis.

\paragraph*{S16 Table.}
\label{si:stats}
{\bf Uncertainty of the main comparisons under three resampling schemes.} Post hoc analysis: images, possible
same-patient clusters, and seeds with clusters.
